\chapterimage{figs/chapterhead.png}
\chapter{Componentes, \textit{model data} e sistema de \textit{plugins}}
\label{chap:componentes_model_data_plugins}

\section{Introdução}
\label{sec:cap8_introducao}

Depois de estudar o núcleo do simulador em seu nível mais central, o próximo passo é compreender de que objetos concretos o modelo é feito. Em outras palavras, se o capítulo anterior respondeu à pergunta ``como o GenESyS representa e executa um modelo?'', este capítulo responde a outra pergunta igualmente importante: ``quais são os elementos que compõem esse modelo dentro do software?''.

No GenESyS, essa resposta não se reduz a um único tipo de classe. O modelo é construído a partir de uma distinção arquitetural fundamental entre:
\begin{itemize}
    \item componentes de modelo;
    \item \textit{model data definitions}.
\end{itemize}

Essa distinção é especialmente importante porque organiza toda a forma como o simulador representa o fluxo principal de um sistema e a infraestrutura de dados que o sustenta. Além disso, é a partir dela que o mecanismo de extensibilidade por \textit{plugins} ganha clareza. Novos componentes e novos dados do modelo podem ser introduzidos ao sistema sem que o núcleo precise ser reescrito como bloco monolítico.

Este capítulo aprofunda essa arquitetura. Primeiro, discute a classe-base \texttt{ModelDataDefinition}. Depois, mostra como \texttt{ModelComponent} a especializa para representar comportamento no fluxo do modelo. Em seguida, discute a importância de dados internos e anexados, a relação entre composição e agregação e, por fim, situa o mecanismo de \textit{plugins} como forma de ampliar o sistema de maneira consistente.

\section{\texttt{ModelDataDefinition} como classe-base do modelo}
\label{sec:cap8_modeldatadefinition}

A classe \texttt{ModelDataDefinition} ocupa papel estrutural no GenESyS. Sua documentação a descreve como a base para qualquer dado do modelo — filas, recursos, variáveis e outros elementos — e também para qualquer componente do modelo. Isso significa que ela define uma infraestrutura comum que será compartilhada tanto por elementos que sustentam a simulação quanto por elementos que participam diretamente do fluxo do modelo.

Essa decisão de projeto é extremamente importante. Em vez de separar completamente os componentes do fluxo e os dados de apoio em hierarquias sem relação, o GenESyS faz com que ambos partilhem uma base comum. Com isso, o software ganha:
\begin{itemize}
    \item um mecanismo uniforme de identificação e nomeação;
    \item suporte comum à persistência;
    \item suporte comum à verificação;
    \item infraestrutura comum para criação de dados internos e anexados;
    \item mecanismos padronizados de tracing e controles de simulação.
\end{itemize}

Em termos arquiteturais, isso torna o modelo mais coerente. Em vez de dois universos sem comunicação, há uma hierarquia unificada na qual componentes e dados são espécies diferentes de elementos do modelo.

\subsection{Responsabilidades centrais de \texttt{ModelDataDefinition}}
\label{subsec:cap8_responsabilidades_modeldatadefinition}

A leitura da classe mostra algumas responsabilidades particularmente relevantes:
\begin{itemize}
    \item manter identidade e nome do elemento;
    \item oferecer infraestrutura de persistência;
    \item permitir verificação local do elemento;
    \item sustentar dados internos e anexados;
    \item oferecer integração com parser e controles de simulação;
    \item prover mecanismos de tracing;
    \item definir ponto de extensão para criação de dados relacionados.
\end{itemize}

Isso significa que \texttt{ModelDataDefinition} não é apenas uma superclasse abstrata conveniente. Ela é um bloco real de infraestrutura sobre o qual o restante do modelo se apoia.

\subsection{Internal data e attached data}
\label{subsec:cap8_internal_attached_data}

Um dos aspectos mais interessantes da classe \texttt{ModelDataDefinition} é a distinção explícita entre \textit{internal data} e \textit{attached data}. Essa distinção, que já foi importante em discussões anteriores do livro, aparece aqui como parte da infraestrutura base do sistema.

\begin{itemize}
    \item \textbf{Internal data} corresponde a elementos internos ao objeto, mantidos por composição.
    \item \textbf{Attached data} corresponde a dados referenciados, mantidos por agregação.
\end{itemize}

Essa distinção é arquiteturalmente muito valiosa. Ela permite que o software represente, de forma clara, tanto elementos que pertencem estruturalmente a um componente quanto elementos que apenas lhe são associados. Em outras palavras, o GenESyS não trata todas as relações entre objetos do modelo da mesma maneira, e isso melhora muito a clareza da semântica interna do sistema.

\subsection{Persistência e verificação como responsabilidades da base}
\label{subsec:cap8_persistencia_verificacao_base}

Outro aspecto importante é que persistência e verificação já aparecem na classe-base. Isso mostra que o sistema espera que elementos do modelo sejam:
\begin{itemize}
    \item carregados;
    \item salvos;
    \item validados;
    \item reconstituídos;
    \item preparados para simulação.
\end{itemize}

Ao concentrar essa infraestrutura na base, o GenESyS evita repetição desnecessária nas subclasses e impõe uma disciplina consistente às extensões do modelo.

\section{\texttt{ModelComponent} como especialização comportamental}
\label{sec:cap8_modelcomponent}

A classe \texttt{ModelComponent} herda de \texttt{ModelDataDefinition} e acrescenta a ela uma dimensão decisiva: comportamento orientado por chegada de entidade e disparo de evento. Sua documentação a descreve como um bloco que representa um comportamento específico a ser simulado e que é ativado quando uma entidade chega ao componente.

Essa formulação é extremamente importante porque define o lugar dos componentes na arquitetura do GenESyS. Um componente não é apenas um objeto visual nem apenas um bloco nominal do modelo. Ele é uma unidade de comportamento dentro do fluxo do sistema modelado.

\subsection{A ideia de fluxo entre componentes}
\label{subsec:cap8_ideia_fluxo_componentes}

Um modelo construído com \texttt{ModelComponent} é, em essência, um conjunto de componentes interconectados. O comportamento global do modelo emerge do percurso das entidades por esse fluxo. Isso explica por que a classe mantém um \texttt{ConnectionManager} e oferece operações de conexão entre componentes.

Essa é uma das grandes diferenças entre \texttt{ModelComponent} e um \textit{model data definition} genérico:
\begin{itemize}
    \item um componente participa diretamente da lógica do fluxo;
    \item um dado do modelo pode sustentar esse fluxo sem ser ele próprio parte do percurso principal.
\end{itemize}

\subsection{Despacho de eventos e \texttt{\_onDispatchEvent}}
\label{subsec:cap8_dispatch_event_ondispatch}

A classe \texttt{ModelComponent} deixa claro, por sua interface, que cada componente concreto deve implementar um comportamento específico de despacho. Isso aparece no método protegido e virtual puro \texttt{\_onDispatchEvent(Entity*, unsigned int)}. É nesse ponto que a semântica própria de cada componente concreto passa a existir.

Essa escolha arquitetural é particularmente boa porque separa:
\begin{itemize}
    \item a infraestrutura comum dos componentes;
    \item o comportamento particular de cada tipo de componente.
\end{itemize}

Com isso, o GenESyS consegue ter uma classe-base forte, mas ainda assim manter subclasses altamente especializadas.

\subsection{Componentes como model data especializados}
\label{subsec:cap8_componentes_modeldata_especializados}

O fato de \texttt{ModelComponent} herdar de \texttt{ModelDataDefinition} tem consequências profundas:
\begin{itemize}
    \item um componente também tem nome, identidade e persistência;
    \item um componente também pode possuir dados internos e anexados;
    \item um componente participa da infraestrutura comum de verificação;
    \item um componente integra-se naturalmente ao restante da semântica do modelo.
\end{itemize}

Isso evita uma separação artificial entre estrutura de dados e estrutura comportamental. Em vez disso, o sistema reconhece que um componente é também um elemento do modelo, apenas dotado de papel comportamental específico.

\section{Fluxo principal e dados de apoio}
\label{sec:cap8_fluxo_principal_dados_apoio}

Uma das ideias mais importantes que o desenvolvedor precisa consolidar é a diferença entre:
\begin{itemize}
    \item o fluxo principal do modelo;
    \item os dados de apoio que sustentam esse fluxo.
\end{itemize}

Na prática:
\begin{itemize}
    \item o fluxo principal é composto por componentes conectados;
    \item os dados de apoio incluem filas, recursos, atributos, variáveis, coletores, contadores e outros objetos relacionados.
\end{itemize}

Essa distinção é central porque orienta tanto a modelagem quanto a extensão do software. Muitos erros de projeto ocorrem justamente quando se tenta tratar um elemento de apoio como se fosse componente do fluxo, ou vice-versa.

\subsection{Por que essa distinção melhora a arquitetura}
\label{subsec:cap8_por_que_distincao_melhora_arquitetura}

Ao separar fluxo principal e dados de apoio, o GenESyS ganha clareza:
\begin{itemize}
    \item o comportamento processual fica concentrado nos componentes;
    \item os objetos de suporte ficam representados como dados do modelo;
    \item a relação entre ambos pode ser expressa de forma explícita e controlada;
    \item o sistema se torna mais extensível.
\end{itemize}

Isso também melhora a legibilidade do software. O desenvolvedor consegue identificar com mais facilidade se um novo elemento que deseja implementar deve ser um componente, um dado do modelo ou uma combinação dos dois por composição ou agregação.

\section{Composição, agregação e o desenho interno do modelo}
\label{sec:cap8_composicao_agregacao}

A infraestrutura de \texttt{ModelDataDefinition} e \texttt{ModelComponent} revela que o GenESyS leva a sério a distinção entre composição e agregação.

Na composição, o elemento interno pertence estruturalmente ao objeto que o contém. Em termos práticos, isso significa que ele é criado, mantido e destruído como parte do ciclo de vida do objeto principal. Na agregação, ao contrário, há apenas referência ou vínculo com um dado existente no modelo, sem absorver seu ciclo de vida.

Essa distinção é particularmente útil quando um componente:
\begin{itemize}
    \item mantém internamente contadores ou estatísticas próprias;
    \item ou referencia objetos externos como filas e recursos.
\end{itemize}

Sem essa separação, o sistema correria o risco de produzir vínculos mal definidos entre elementos do modelo, dificultando persistência, remoção, verificação e reuso.

\subsection{Internal data como materialização de composição}
\label{subsec:cap8_internaldata_composicao}

Quando um componente ou dado do modelo mantém \textit{internal data}, ele está explicitando uma relação de composição. Isso costuma ser adequado para elementos que:
\begin{itemize}
    \item existem apenas para apoiar o funcionamento interno daquele objeto;
    \item não fazem sentido como objetos independentes do modelo;
    \item devem desaparecer junto com o objeto principal.
\end{itemize}

\subsection{Attached data como materialização de agregação}
\label{subsec:cap8_attacheddata_agregacao}

Quando um objeto mantém \textit{attached data}, o que há é uma relação de agregação. O dado referenciado continua existindo no modelo como objeto próprio, podendo ser compartilhado, referenciado por outros elementos e persistido independentemente.

Essa separação entre composição e agregação torna a arquitetura do GenESyS muito mais limpa e semanticamente mais expressiva.

\section{O sistema de \textit{plugins}}
\label{sec:cap8_sistema_plugins}

A presença de \textit{plugins} é uma das características mais importantes do GenESyS como software extensível. A própria classe \texttt{Simulator} já deixa isso evidente ao expor um \texttt{PluginManager}. Isso mostra que o mecanismo de \textit{plugins} não é improvisado nem periférico; ele é parte estrutural do sistema.

Para o desenvolvedor, a consequência é clara: ampliar o GenESyS não significa necessariamente modificar diretamente o kernel. Em muitos casos, a estratégia mais adequada é introduzir novas capacidades por meio de \textit{plugins}. Isso vale especialmente para:
\begin{itemize}
    \item novos componentes;
    \item novos \textit{data definitions};
    \item novas extensões de linguagem;
    \item novos comportamentos especializados.
\end{itemize}

\subsection{Por que \textit{plugins} são importantes}
\label{subsec:cap8_por_que_plugins_importantes}

Sem \textit{plugins}, o crescimento do sistema dependeria de alterações cada vez mais profundas no núcleo. Isso tenderia a tornar o software mais rígido e mais difícil de manter. Com \textit{plugins}, o projeto ganha:
\begin{itemize}
    \item modularidade;
    \item extensibilidade;
    \item especialização por domínio;
    \item maior separação entre infraestrutura estável e funcionalidades adicionais.
\end{itemize}

Essa é uma decisão arquitetural especialmente apropriada para um simulador que pretende ser genérico e expansível.

\subsection{Componentes e dados como extensões naturais}
\label{subsec:cap8_componentes_dados_extensoes_naturais}

Do ponto de vista do modelo, muitos \textit{plugins} se materializam justamente como novos componentes e novos dados. Isso significa que a arquitetura discutida até aqui — \texttt{ModelDataDefinition}, \texttt{ModelComponent}, internal data, attached data — não é apenas uma questão de elegância interna. Ela é o fundamento sobre o qual a extensibilidade do sistema se apoia.

Em outras palavras, compreender a arquitetura dos elementos do modelo é também compreender a arquitetura dos \textit{plugins}.

\section{Como pensar uma nova extensão}
\label{sec:cap8_como_pensar_nova_extensao}

Um desenvolvedor que pretenda ampliar o GenESyS deve começar por uma pergunta muito simples:
\begin{quote}
    O que exatamente estou tentando introduzir no sistema?
\end{quote}

A partir dessa pergunta, costuma ser possível distinguir alguns casos:

\begin{itemize}
    \item um novo comportamento no fluxo do modelo;
    \item um novo dado de apoio ao modelo;
    \item uma nova combinação entre comportamento e dados internos;
    \item uma nova necessidade de linguagem ou parser;
    \item uma nova aplicação ou ferramenta.
\end{itemize}

Se a necessidade for um novo comportamento no fluxo, a extensão provavelmente terá a forma de novo \texttt{ModelComponent}. Se for um novo objeto de apoio, provavelmente tomará a forma de novo \texttt{ModelDataDefinition}. Se envolver ambos, será necessário pensar cuidadosamente em composição, agregação e relações internas do novo elemento.

\section{Leitura errada que este capítulo evita}
\label{sec:cap8_leitura_errada_evita}

Este capítulo procura evitar uma leitura errada muito comum: imaginar que o modelo do GenESyS seja composto apenas por caixas visuais na GUI. Na realidade, o modelo é uma estrutura muito mais rica, sustentada por uma hierarquia de classes e por uma distinção arquitetural clara entre comportamento e dados.

Também procura evitar outro erro: imaginar que \textit{plugins} sejam apenas detalhes de carregamento dinâmico. No GenESyS, eles são parte da lógica de crescimento do sistema. A extensibilidade não é apêndice; é princípio arquitetural.

\begin{table}[htbp]
    \centering
    \caption{Distinções centrais para compreender a estrutura do modelo no GenESyS.}
    \label{tab:cap8_distincoes_centrais}
    \small
    \begin{tabular}{|p{3.8cm}|p{7.4cm}|}
        \hline
        \textbf{Distinção} & \textbf{Significado arquitetural} \\
        \hline
        \texttt{ModelComponent} vs. \texttt{ModelDataDefinition} & Diferença entre comportamento no fluxo e dados/infraestrutura do modelo \\
        \hline
        Fluxo principal vs. dados de apoio & Diferença entre percurso processual e elementos que sustentam a execução \\
        \hline
        Internal data vs. attached data & Diferença entre composição e agregação dentro do modelo \\
        \hline
        Kernel vs. \textit{plugins} & Diferença entre infraestrutura central e extensões especializadas \\
        \hline
    \end{tabular}
\end{table}

\section{Considerações Finais}
\label{sec:cap8_consideracoes_finais}

Este capítulo mostrou que o modelo do GenESyS é estruturado a partir de uma arquitetura coerente e extensível. A classe \texttt{ModelDataDefinition} fornece a infraestrutura base para elementos do modelo. A classe \texttt{ModelComponent} especializa essa base para representar comportamento no fluxo. A distinção entre fluxo principal e dados de apoio, bem como entre composição e agregação, organiza de forma clara a semântica interna do sistema. Finalmente, o mecanismo de \textit{plugins} aparece como consequência natural dessa arquitetura, permitindo que o simulador seja ampliado sem perder coerência.

Com isso, o leitor já dispõe da base necessária para compreender como novos elementos podem ser introduzidos no sistema. O próximo passo natural é estudar o subsistema que conecta linguagem, expressões e semântica do modelo: o parser do GenESyS e sua relação com a linguagem interna do simulador.

Teste seu entendimento desse conteúdo respondendo as seguintes perguntas:

\begin{enumerate}
    \item Qual é a vantagem arquitetural de fazer \texttt{ModelComponent} herdar de \texttt{ModelDataDefinition}?

    \item Como a distinção entre fluxo principal e dados de apoio ajuda a compreender a estrutura do modelo?

    \item Qual é a diferença entre \textit{internal data} e \textit{attached data} no desenho interno do GenESyS?

    \item Por que o sistema de \textit{plugins} deve ser entendido como parte estrutural da arquitetura e não apenas como mecanismo opcional de extensão?
\end{enumerate}